\documentclass[11pt]{article}
\usepackage{mathblatt}
% Präambel des Setzers (werkzeuge/setzer-praeambel.tex, wortgleich); dazu \gteil (Muster T2U)
\makeatletter
\setlength{\mbpfbe}{0mm}% keine Punkte (pruefung.md § 5: nur beim Prüfstein)
\newlength{\gEin}\setlength{\gEin}{0pt}
\let\gSchrift\relax
\renewcommand{\pfkreuzzeile}[1]{\par\noindent\hangindent3.2em\hangafter1\hspace*{1.6em}$\square$\ #1\par}
\newcommand{\gkreuz}[1]{$\square$\ #1\hspace{2em}}
\newcommand{\gkreuzzeile}[2]{\par\noindent\hangindent2.2em\hangafter1\mbox{#1\,$\square$}\ #2\par}% Buchstabe, Option: eine je Zeile (neben Skizze)
% links, was man ansieht, rechts, was man tut; Spaltengrenze überall an derselben Stelle (gemeinsam.md „Satz“)
% v0.17: Fußeintrag „Nr: Wert“ – kurz in einer Box (nie getrennt), lang umbrechbar mit der Nummer am Anfang
% (vorher \mbox: lange Lösungen liefen über den Rand, AS40 138 pt); robust, weil \fusshilfe in die .aux schreibt
\NewDocumentCommand{\gfuss}{m m}{\sbox0{#1:~#2}\ifdim\wd0<0.45\textwidth\mbox{#1:~#2}\else\mbox{#1:}\nobreak\ #2\fi}
\newsavebox{\gbox}\newlength{\gLinks}
% v0.10 (Lehrer 13.10.): keine starre Halbierung – die linke Spalte ist so breit wie die Figur (+ 6 mm), die
% Antwort steht direkt daneben; ist die Figur breiter als 0,62 Zeile, steht die Antwort darunter
\newcommand{\gzwei}[2]{\par\smallskip\noindent\sbox{\gbox}{#1}%
  \setlength{\gLinks}{\dimexpr\wd\gbox+6mm\relax}%
  \ifdim\gLinks>0.62\linewidth
    \usebox{\gbox}\par\smallskip\noindent\begin{minipage}{\linewidth}#2\end{minipage}%
  \else
    \begin{minipage}[t]{\gLinks}\vspace{0pt}\usebox{\gbox}\end{minipage}%
    \begin{minipage}[t]{\dimexpr\linewidth-\gLinks\relax}\vspace{0pt}\raggedright #2\end{minipage}%
  \fi\par}
% v0.9 (12.10., pruefung.md § 5 „Marke mit B, Bedingungen rechts“): drittes Argument von pfaufg = Bedingungen
% der Prüfung (v0.10: \pfohnetr = TR durchgestrichen, „FOR“), klein grau im rechten Rand auf der Grundlinie der Nummer, Zeilen mit \\
\newcommand{\gbedrand}[2]{\ifblank{#2}{}{\rlap{\hspace*{\dimexpr#1+2mm\relax}%
  \parbox[t]{15mm}{\raggedright\scriptsize\leavevmode\color{mbgrau}\strut#2}}}}
% Teil eines Bündels (a, b …): Marke links im Rand, Bedingungen rechts im Rand, auf der Zeile des Teils
\newcommand{\gteilrand}[2]{\llap{\mbpfrandmarke{#1}\hspace*{\mbpfnr}}\gbedrand{\linewidth}{#2}}
% Teil eines Bündels: Buchstabe in eigener Spalte, Inhalt eingerückt (\gEin = Einzug, damit die Spaltengrenze von
% \gzwei an derselben Stelle bleibt)
\newenvironment{gbteil}[3]{\par\smallskip\noindent\gteilrand{#2}{#3}\makebox[\mbpfteil][l]{\textbf{#1}}%
  \begin{minipage}[t]{\dimexpr\linewidth-\mbpfteil\relax}\setlength{\gEin}{\mbpfteil}\raggedright\strut\ignorespaces}%
 {\end{minipage}\par\addvspace{6pt}}% v0.18 (PG80 Nr. 1): Abstand, damit die Linie von a) nicht an b) rutscht
% Grafik höchstens so breit wie die Spalte (Bündel nebeneinander, Koordinatensystem ohne Karo)
\newsavebox{\gpbox}
\newcommand{\gpassend}[1]{\sbox{\gpbox}{#1}\ifdim\wd\gpbox>\linewidth\typeout{GPASSEND: verkleinert \the\wd\gpbox\space auf \the\linewidth}\resizebox{\linewidth}{!}{\usebox{\gpbox}}\else\usebox{\gpbox}\fi}
% Skizzenfeld (v0.10, Lehrer 13.10.): Karo ab dem linken Rand über die Zeile, „Skizze“ klein grau oben links im Feld
% (wie \pfskizzenfeld der Vorlage); der rechte Rand bleibt den Bedingungen
\newcommand{\gskizze}[1]{\par\smallskip\noindent\begin{tikzpicture}
  \clip (0,0) rectangle ({\linewidth-1mm},{#1*10mm});
  \draw[black!16,line width=0.3pt,step=5mm] (0,0) grid ({\linewidth},{#1*10mm});
  \node[anchor=north west,inner sep=2pt,font=\scriptsize,text=mbgrau,fill=white] at (0,{#1*10mm}) {Skizze};
  \end{tikzpicture}\par}
\RenewDocumentEnvironment{pfaufg}{m m m}{%
  \begin{lrbox}{\mbaufgabebox}\begin{minipage}[t]{\linewidth}%
  \gSchrift\noindent\hspace*{\gEin}\mbpfrandmarke{#1}\gbedrand{\linewidth-\gEin}{#3}%
  \begin{minipage}[t]{\mbpfnr}\strut\textbf{#2}\end{minipage}%
  \begin{minipage}[t]{\dimexpr\linewidth-\gEin-\mbpfnr\relax}\raggedright\strut\ignorespaces}%
 {\end{minipage}%
  \end{minipage}\end{lrbox}%
  \setlength{\mbaufgabehoehe}{\dimexpr\ht\mbaufgabebox+\dp\mbaufgabebox\relax}%
  \ifdefempty{\gkopfnach}{\par\addvspace{5pt}\Needspace*{\mbaufgabehoehe}}{%
    % Stufenkopf, der mit dieser (ersten) Aufgabe zusammenbleibt: Platz für beide verlangen, dann den Kopf setzen
    % (12.10., Quadratische: ein Kopf stand allein unten, die Aufgabe mit Koordinatensystem auf der nächsten Seite)
    % v0.17 (AS40 „Terme deuten“ allein unten): hinter dem Kopf steht das Sprungziel (\special, ein Knoten, nach
    % dem die Seite brechen durfte) – darum \nobreak nach dem Kopf; Vorhalt 56 pt = Kopf mit Linie und Abständen
    \par\Needspace*{\dimexpr\mbaufgabehoehe+56pt\relax}\gkopfnach\par\nobreak\gdef\gkopfnach{}\addvspace{5pt}}%
  \noindent\usebox{\mbaufgabebox}\par\addvspace{5pt}}
% v0.18 (plan.md § 10, 15.10. l): Abschnittsbuchstabe A, B, C vor dem Namen im Kopf – als Vorsatz, damit der
% Name im TeX gleich bleibt (Bestellkarte, Sprungziele lesen ihn)
\def\gbstakt{}\newcommand{\gbstvor}[1]{\gdef\gbstakt{#1}}
\NewCommandCopy\gpfstufekopfalt\pfstufekopf
\RenewDocumentCommand{\pfstufekopf}{o m m}{\IfValueTF{#1}{\gpfstufekopfalt[#1]}{\gpfstufekopfalt}%
  {\ifx\gbstakt\empty\else\gbstakt\enspace\fi #2}{#3}\gdef\gbstakt{}}
% v0.18 (16.10. e): Begriff grau, ohne Nummer, ohne Rahmen über der ersten Aufgabe der Stufe
\newcommand{\gbegriff}[1]{\par\nobreak\noindent\hspace*{\mbpfnr}\parbox[t]{\dimexpr\linewidth-\mbpfnr\relax}{\raggedright\small\color{mbgrau}#1}\par\nobreak\addvspace{3pt}}
% v0.18 (16.10. g): Lösungsblatt – Kopf mit kräftigem Strich, Abschnitt klein grau als Zwischenzeile
\newcommand{\glkopf}[1]{\noindent{\large\bfseries Lösungen \textperiodcentered{} #1}\par\vspace{3pt}\hrule height 1.2pt\par\medskip}
\newcommand{\glzwischen}[1]{\par\addvspace{8pt}\Needspace*{4\baselineskip}\noindent{\footnotesize\color{mbgrau}#1}\par\nobreak\vspace{1pt}}
% v0.18 (15.10. o, w): Blattende – Kern ohne Rahmen, dann graue Zeile „Mehr: n Aufgaben“
\newcommand{\gkern}[1]{\par\addvspace{14pt}\Needspace*{6\baselineskip}\noindent{\bfseries Kern}\par\nobreak\smallskip
  {\small\raggedright #1\par}}
\newcommand{\gmehr}[1]{\par\addvspace{10pt}\noindent{\footnotesize\color{mbgrau}Mehr: #1 Aufgaben}\par}
% Stufenkopf vormerken: \gkopfvor{<Kopf>} – die nächste pfaufg setzt ihn, sobald beide zusammen Platz haben
\def\gkopfnach{}
\newcommand{\gkopfvor}[1]{\gdef\gkopfnach{#1}}
% Teilaufgabe mit Skizze links in gzwei: Buchstabe in eigener Spalte, darunter Text und Skizze
\newcommand{\gteil}[2]{\makebox[\mbpfteil][l]{\textbf{#1}}\parbox[t]{\dimexpr0.5\textwidth-\gEin-\mbpfnr-5mm-\mbpfteil\relax}{\raggedright #2}}
% Stufenkopf der Vorlage (mathblatt.sty 2026-10-12: Linie, Abstand, 13 pt, \pfjahre) ohne ihren Vorhalt von
% 14 Zeilen: der Setzer fordert vor dem Kopf 0,16 Seitenhöhe (Kopf + erste Aufgabe), mehr nicht
\setlength{\mbpfstufevorhalt}{0pt}
\makeatother
\newcommand{\karo}[1]{\par\smallskip\noindent\begin{tikzpicture}
  \clip (0,0) rectangle ({\linewidth-1mm},{#1*10mm});
  \draw[black!16,line width=0.3pt,step=5mm] (0,0) grid ({\linewidth},{#1*10mm});
  \end{tikzpicture}\par}
\newcommand{\kk}[1]{$\square$\,#1\hspace{0.9em}}
\newcommand{\linie}[1]{\,{\color{black!45}\rule[-1pt]{#1}{0.4pt}}\,}
\newcommand{\pkt}{\textperiodcentered{}}
% v0.14 (14.10.): Stufenkopf mit fünf Punkten (Option = Jahre für die Bestellkarte) oder „zuletzt 19“
\NewDocumentCommand{\gpunkte}{o m}{#2}
\newcommand{\gzuletzt}[1]{zuletzt #1}
% v0.7 (12.10.): Schülerfrage unter dem Stufenkopf; Ankreuzoption als Bild (Buchstabe und Kästchen darunter)
\newcommand{\gfrage}[1]{\par\nobreak\noindent\hspace*{\mbpfnr}{\small\itshape #1}\par\nobreak\addvspace{2pt}}
\newlength{\gbildbreite}\setlength{\gbildbreite}{0.32\linewidth}% Setzer setzt sie je Reihe (vier Bilder: 0,24)
\newcommand{\gbildoption}[3]{\begin{minipage}[t][#3][b]{\gbildbreite}\centering #2\par\smallskip\mbox{#1\,$\square$}\end{minipage}}
% Lineare (12.10.): Bild mit Buchstabe ohne Kästchen (Diagramm zuordnen); Aussage mit „wahr / falsch“ zum Ankreuzen
\newcommand{\gbildname}[3]{\begin{minipage}[t][#3][b]{\gbildbreite}\centering #2\par\smallskip\mbox{\textbf{#1}}\end{minipage}}
\newcommand{\gwahrfalsch}[2]{\par\noindent\hangindent2.2em\hangafter1\makebox[2.2em][l]{(#1)}#2\hfill\mbox{wahr $\square$\quad falsch $\square$}\par}
\begin{document}

\pfheftstil
\pflaufkopf[P10 \pkt{} Prüfung]{Volumen und Oberfläche}{PG90-7}\pfseitevon
\pfheftkopf{Volumen und Oberfläche}{}
\pfstufenzeile{A~Volumenformel aufstellen \pkt{} B~Volumen gesucht \pkt{} C~Volumenformel umstellen \pkt{} D~Oberfläche gesucht \pkt{} E~Körper vergleichen}

% ===================== Stufe: Volumenformel aufstellen =====================
\Needspace*{0.16\textheight}\gbstvor{A}\pfstufekopf{Volumenformel aufstellen}{\gpunkte[]{\textcolor{black!25}{$\circ$}\,\textcolor{black!25}{$\circ$}\,\textcolor{black!25}{$\circ$}\,\textcolor{black!25}{$\circ$}\,\textcolor{black!25}{$\circ$}}}\special{pdf:dest (stufe-volumenformel-aufstellen) [@thispage /XYZ @xpos @ypos null]}\gfrage{Such in deiner Formelsammlung das Volumen des Zylinders: Seite~\linie{12mm}}
% Schülerfrage der Stufe (Gliederung; Einstieg nur von dort, nie aus der Bank – 12.10.)
\begin{pfaufg}{}{1.}{}
Wie viel passt hinein? Schreibe nur die Formel auf, rechne nicht.
\gzwei{\begin{tikzpicture}[line width=0.7pt,font=\small]\draw (0,1.6) ellipse (0.5 and 0.17);\draw (-0.5,0)--(-0.5,1.6);\draw (0.5,0)--(0.5,1.6);\draw (-0.5,0) arc (180:360:0.5 and 0.17);\draw[dashed,line width=0.5pt] (-0.5,0) arc (180:0:0.5 and 0.17);\node at (0,-0.55) {a)};\draw (1.2,0)--(2.4,0)--(1.8,1.0)--cycle;\draw (2.4,0)--(3.1,0.5)--(2.5,1.5)--(1.8,1.0);\draw[dashed,line width=0.5pt] (1.2,0)--(1.9,0.5)--(3.1,0.5);\draw[dashed,line width=0.5pt] (1.9,0.5)--(2.5,1.5);\node at (2.1,-0.55) {b)};\draw (4.2,1.6) ellipse (0.5 and 0.17);\draw (3.7,1.6)--(4.2,0)--(4.7,1.6);\node at (4.2,-0.55) {c)};\end{tikzpicture}}{\vspace{-\baselineskip}\par\noindent\makebox[\mbpfteil][l]{\textbf{a)}}Konservendose, 7 cm hoch, Radius 3,5 cm\hfill\linie{30mm}\par\noindent\makebox[\mbpfteil][l]{\textbf{b)}}Zelt, 4 m lang, Vorderseite ein Dreieck\hfill\linie{30mm}\par\noindent\makebox[\mbpfteil][l]{\textbf{c)}}Eistüte, Radius 3 cm, Höhe 12 cm\hfill\linie{30mm}}
\par\fusshilfe{\gfuss{1}{a) Zylinder; b) Prisma; c) Kegel}}
\end{pfaufg}
% SETZER: fehlt – keine echte Aufgabe mit auf_blatt = ja in msa/mehr-p10/*.csv

% ===================== Stufe: Volumen direkt =====================
\gkopfvor{\gbstvor{B}\pfstufekopf{Volumen gesucht}{\gpunkte[26 24 23 22]{\textcolor{black!80}{$\bullet$}\,\textcolor{black!25}{$\circ$}\,\textcolor{black!80}{$\bullet$}\,\textcolor{black!80}{$\bullet$}\,\textcolor{black!80}{$\bullet$}}}\special{pdf:dest (stufe-volumen-gesucht) [@thispage /XYZ @xpos @ypos null]}}
% kein Einstieg: keine Schülerfrage in der Gliederung (Schülerfrage, -Fälle, -Form, -Weg, -Skizze; oder Frage) – Einstieg nur von dort, nie aus der Bank (pruefung.md § 5, 12.10.)
% 2023-OS-K5b-m1 (Original 2023-OS-K5b, Schwere I, rang 3; herausgelöste Fassung: keine)
\begin{pfaufg}{nach~P10\\’23~\pkt{}~5b}{2.}{}
Eine Konservendose hat die Form eines Zylinders mit dem Radius 3,5 cm und der Höhe 7 cm. Auf der Dose steht: Inhalt etwa 270 ml. Weise nach, dass das stimmt (1 cm³ = 1 ml).
\gzwei{\zylinder{0.8}{1.9}{}{}}{\karo{2}}
\fusshilfe{\gfuss{2}{269,4 cm³}}
\end{pfaufg}
% 2024-OS-K4c-m1 (Original 2024-OS-K4c, Schwere II, rang 6; herausgelöste Fassung: 2024-OS-K4c-h1)
\begin{pfaufg}{nach~P10\\’24~\pkt{}~4c}{3.}{FOR}
Eine Firma stellt Kegel mit dem Radius r = 25 cm und der Höhe h = 60 cm her. Eine andere Firma verschickt denselben Kegel in einer zylinderförmigen Verpackung; innen ist sie 61 cm hoch und hat den Radius 25,5 cm. Der freie Raum wird mit Füllmaterial ausgefüllt. Ermittle, wie viele cm³ Füllmaterial man braucht.
\gskizze{3}
\fusshilfe{\gfuss{3}{85342 cm³}}
\end{pfaufg}

% ===================== Stufe: rückwärts: Radius oder Höhe aus Volumen =====================
\Needspace*{0.16\textheight}\gbstvor{C}\pfstufekopf{Volumenformel umstellen}{\gpunkte[23 22]{\textcolor{black!25}{$\circ$}\,\textcolor{black!25}{$\circ$}\,\textcolor{black!25}{$\circ$}\,\textcolor{black!80}{$\bullet$}\,\textcolor{black!80}{$\bullet$}}}\special{pdf:dest (stufe-volumenformel-umstellen) [@thispage /XYZ @xpos @ypos null]}
% Schülerfrage der Stufe (Gliederung; Einstieg nur von dort, nie aus der Bank – 12.10.)
\begin{pfaufg}{}{4.}{}
Eine Säule hat eine quadratische Grundfläche. Stelle die Formel V = a² · h nach a um.
\gzwei{\quader{0.7}{0.7}{1.6}{a}{a}{h}}{\vspace{-\baselineskip}\par\noindent \hfill\linie{30mm}}
\par\fusshilfe{\gfuss{4}{$a = \sqrt{\tfrac{V}{h}}$}}
\end{pfaufg}
% 2023-OS-K5d-m1 (Original 2023-OS-K5d, Schwere II, rang 7; herausgelöste Fassung: keine)
\begin{pfaufg}{nach~P10\\’23~\pkt{}~5d}{5.}{FOR}
Eine Konservendose hat die Form eines Zylinders mit dem Radius 3,5 cm und der Höhe 7 cm. Mit denselben Deckeln (r = 3,5 cm) soll auch eine größere Dose mit 1 Liter Inhalt gebaut werden. Ermittle ihre Höhe.
\gzwei{\zylinder{0.8}{1.9}{}{}}{\karo{2}}
\fusshilfe{\gfuss{5}{26,0 cm}}
\end{pfaufg}
% 2022-OS-K2d-m1 (Original 2022-OS-K2d, Schwere II, rang 8; herausgelöste Fassung: keine)
\begin{pfaufg}{nach~P10\\’22~\pkt{}~2d}{6.}{FOR}
Ein Becher hat die Form eines Zylinders ohne Deckel. Er ist h = 7,5 cm hoch und hat den Radius r = 3,4 cm. Ein größerer Becher soll 500 ml fassen und wieder 7,5 cm hoch sein. Berechne seinen Radius.
\gzwei{\begin{tikzpicture}[baseline=(current bounding box.north),line width=0.7pt]\fill[black!12] (0,2) ellipse (0.8 and 0.28);\draw (0,2) ellipse (0.8 and 0.28);\draw (-0.8,0) -- (-0.8,2);\draw (0.8,0) -- (0.8,2);\draw (-0.8,0) arc (180:360:0.8 and 0.28);\draw[dashed,line width=0.5pt] (-0.8,0) arc (180:0:0.8 and 0.28);\end{tikzpicture}}{\karo{2}}
\fusshilfe{\gfuss{6}{4,6 cm}}
\end{pfaufg}

% ===================== Stufe: Mantelfläche mit Kosten =====================
\Needspace*{0.16\textheight}\gbstvor{D}\pfstufekopf{Oberfläche gesucht}{\gpunkte[26]{\textcolor{black!80}{$\bullet$}\,\textcolor{black!25}{$\circ$}\,\textcolor{black!25}{$\circ$}\,\textcolor{black!25}{$\circ$}\,\textcolor{black!25}{$\circ$}}}\special{pdf:dest (stufe-oberflache-gesucht) [@thispage /XYZ @xpos @ypos null]}
% Schülerfrage der Stufe (Gliederung; Einstieg nur von dort, nie aus der Bank – 12.10.)
\begin{pfaufg}{}{7.}{}
Du schneidest den Mantel eines Zylinders gerade auf und rollst ihn flach aus. Welche Figur entsteht, und wie lang sind ihre Seiten?
\gzwei{\zylinder{0.6}{1.5}{r}{h}}{\vspace{-\baselineskip}\par\noindent \hfill\linie{30mm}}
\par\fusshilfe{\gfuss{7}{Rechteck, 2 · $\pi$ · r lang und h hoch}}
\end{pfaufg}
% 2026-FOR-K2b-m1 (Original 2026-FOR-K2b, Schwere II, rang 5; herausgelöste Fassung: keine)
\begin{pfaufg}{nach~P10\\’26~\pkt{}~2b}{8.}{}
Ein Turm besteht unten aus einem Zylinder, sein Dach ist ein Kegel. Bekannt sind: h = 22,0 m (Höhe des Zylinders), r = 4,2 m (Radius von Zylinder und Kegel) und s = 7,9 m (Mantellinie des Kegels). Das Dach wird mit Ziegeln gedeckt; 1 m² Ziegel kostet 24,00 €. Berechne die Kosten für die Ziegel.
\gzwei{\begin{tikzpicture}[baseline=(current bounding box.north),line width=0.7pt]\draw (-0.9,0) arc (180:360:0.9 and 0.315);\draw[dashed,line width=0.5pt] (-0.9,0) arc (180:0:0.9 and 0.315);\draw (-0.9,0) -- (-0.9,2.4);\draw (0.9,0) -- (0.9,2.4);\draw (0,2.4) ellipse (0.9 and 0.315);\draw (-0.9,2.4) -- (0,3.9) -- (0.9,2.4);\node[right,font=\small] at (0.95,1.2) {$h$};\fill (0,0) circle (0.8pt);\draw[dashed,line width=0.5pt] (0,0) -- (0.9,0) node[midway,fill=white,inner sep=1pt,font=\small] {$r$};\path (0.9,2.4) -- (0,3.9) node[midway,above right,inner sep=1.5pt,font=\small] {$s$};\end{tikzpicture}}{\karo{3}}
\fusshilfe{\gfuss{8}{104,24 m²; 2501,71 €}}
\end{pfaufg}

% ===================== Stufe: vergleichen und urteilen =====================
\Needspace*{0.16\textheight}\gbstvor{E}\pfstufekopf{Körper vergleichen}{\gpunkte[26]{\textcolor{black!80}{$\bullet$}\,\textcolor{black!25}{$\circ$}\,\textcolor{black!25}{$\circ$}\,\textcolor{black!25}{$\circ$}\,\textcolor{black!25}{$\circ$}}}\special{pdf:dest (stufe-korper-vergleichen) [@thispage /XYZ @xpos @ypos null]}
% Schülerfrage der Stufe (Gliederung; Einstieg nur von dort, nie aus der Bank – 12.10.)
\begin{pfaufg}{}{9.}{}
Wie oft so groß wird das Volumen? Schätze, ohne zu rechnen.
\par\smallskip
\par\noindent\makebox[\mbpfteil][l]{\textbf{a)}}Zylinder, die Höhe wird verdoppelt\hfill\linie{30mm} -mal so groß
\par\noindent\makebox[\mbpfteil][l]{\textbf{b)}}Zylinder, der Radius wird verdoppelt\hfill\linie{30mm} -mal so groß
\par\noindent\makebox[\mbpfteil][l]{\textbf{c)}}Zylinder statt Kegel, Radius und Höhe gleich\hfill\linie{30mm} -mal so groß
\par\fusshilfe{\gfuss{9}{a) 2-mal so groß; b) 4-mal so groß; c) 3-mal so groß}}
\end{pfaufg}
% 2026-FOR-K2d-m1 (Original 2026-FOR-K2d, Schwere III, rang 9; herausgelöste Fassung: keine)
\begin{pfaufg}{nach~P10\\’26~\pkt{}~2d}{10.}{FOR}
Ein Turm besteht unten aus einem Zylinder, sein Dach ist ein Kegel. Bekannt sind: h = 22,0 m (Höhe des Zylinders), r = 4,2 m (Radius von Zylinder und Kegel) und s = 7,9 m (Mantellinie des Kegels). Ein Zylinder und ein Kegel sind gleich hoch. Jemand behauptet: „Ist der Radius des Kegels doppelt so groß wie der des Zylinders, haben beide dasselbe Volumen.“ Untersuche mit einer Rechnung, ob das stimmt.
\gskizze{3}
\fusshilfe{\gfuss{10}{falsch}}
\end{pfaufg}

\gmehr{8}
\end{document}